% GENERATED FILE. Edit canonical lessons, then run npm run generate:books.
% canonical-source-hash: fnv1a64:26a5657a1b217d34
% canonical-lessons: GE-C164-nachbarin, GE-C164-cousine, GE-C164-enkelin, GE-C164-baby, GE-C164-mensch

\chapter{A neighbour, A cousin, A granddaughter, A baby, A human being}
\label{ch:ge-a-neighbour-a-cousin-a-granddaughter-a-baby-a-human-being}
\hlchaptermodality{164}

\begin{chapteropening}
\textbf{By the end of this chapter:} \emph{I can say a neighbour, a cousin, a granddaughter, a baby and a human being.}

Complete the last lesson of chapter 164: i can say a neighbour, a cousin, a granddaughter, a baby and a human being.
\end{chapteropening}

\section[die Nachbarin]{die Nachbarin — a neighbour (woman)}

\label{lesson:GE-C164-nachbarin}



\begin{quote}
Before the new one: say the German for a party, then the German for an invitation.
\end{quote}

\subsection*{You'll want to know: die Nachbarin}
\textbf{die Nachbarin} — \textquotedblleft{}a neighbour (woman)\textquotedblright{}.

\begin{grammarlens}[title={What you've built}]
One more word for this chapter.
\end{grammarlens}

\subsection*{Your turn}
\begin{itemize}
\raggedright
  \item \emph{Say it:} \emph{die Nachbarin}
  \item \emph{Say it:} \emph{die Nachbarin}, once more
  \item \emph{Say it:} say \emph{die Nachbarin}
  \item \emph{Recall:} say the German for a party, then the German for an invitation, then say \emph{die Nachbarin} again
\end{itemize}

\begin{tcolorbox}[breakable,colback=teal!4,colframe=teal!35!black,title={Before you move on}]
What does die Nachbarin mean? (\textquotedblleft{}A neighbour (woman)\textquotedblright{}.) Say it once more.
\end{tcolorbox}
\section[die Cousine]{die Cousine — a cousin (woman)}

\label{lesson:GE-C164-cousine}



\begin{quote}
Before the new one: say the German for an invitation, then the German for a neighbour.
\end{quote}

\subsection*{You'll want to know: die Cousine}
\textbf{die Cousine} — \textquotedblleft{}a cousin (woman)\textquotedblright{}.

\begin{grammarlens}[title={What you've built}]
One more word for this chapter.
\end{grammarlens}

\subsection*{Your turn}
\begin{itemize}
\raggedright
  \item \emph{Say it:} \emph{die Cousine}
  \item \emph{Say it:} \emph{die Cousine}, once more
  \item \emph{Say it:} say \emph{die Cousine}
  \item \emph{Recall:} say the German for an invitation, then the German for a neighbour, then say \emph{die Cousine} again
\end{itemize}

\begin{tcolorbox}[breakable,colback=teal!4,colframe=teal!35!black,title={Before you move on}]
What does die Cousine mean? (\textquotedblleft{}A cousin (woman)\textquotedblright{}.) Say it once more.
\end{tcolorbox}
\section[die Enkelin]{die Enkelin — a granddaughter}

\label{lesson:GE-C164-enkelin}



\begin{quote}
Before the new one: say the German for a neighbour, then the German for a cousin.
\end{quote}

\subsection*{You'll want to know: die Enkelin}
\textbf{die Enkelin} — \textquotedblleft{}a granddaughter\textquotedblright{}.

\begin{grammarlens}[title={What you've built}]
One more word for this chapter.
\end{grammarlens}

\subsection*{Your turn}
\begin{itemize}
\raggedright
  \item \emph{Say it:} \emph{die Enkelin}
  \item \emph{Say it:} \emph{die Enkelin}, once more
  \item \emph{Say it:} say \emph{die Enkelin}
  \item \emph{Recall:} say the German for a neighbour, then the German for a cousin, then say \emph{die Enkelin} again
\end{itemize}

\begin{tcolorbox}[breakable,colback=teal!4,colframe=teal!35!black,title={Before you move on}]
What does die Enkelin mean? (\textquotedblleft{}A granddaughter\textquotedblright{}.) Say it once more.
\end{tcolorbox}
\section[das Baby]{das Baby — a baby}

\label{lesson:GE-C164-baby}



\begin{quote}
Before the new one: say the German for a cousin, then the German for a granddaughter.
\end{quote}

\subsection*{You'll want to know: das Baby}
\textbf{das Baby} — \textquotedblleft{}a baby\textquotedblright{}.

\begin{grammarlens}[title={What you've built}]
One more word for this chapter.
\end{grammarlens}

\subsection*{Your turn}
\begin{itemize}
\raggedright
  \item \emph{Say it:} \emph{das Baby}
  \item \emph{Say it:} \emph{das Baby}, once more
  \item \emph{Say it:} say \emph{das Baby}
  \item \emph{Recall:} say the German for a cousin, then the German for a granddaughter, then say \emph{das Baby} again
\end{itemize}

\begin{tcolorbox}[breakable,colback=teal!4,colframe=teal!35!black,title={Before you move on}]
What does das Baby mean? (\textquotedblleft{}A baby\textquotedblright{}.) Say it once more.
\end{tcolorbox}
\section[der Mensch]{der Mensch — a human being}

\label{lesson:GE-C164-mensch}



\begin{quote}
Before the new one: say the German for a granddaughter, then the German for a baby.
\end{quote}

\subsection*{You'll want to know: der Mensch}
\textbf{der Mensch} — \textquotedblleft{}a human being\textquotedblright{}.

\begin{grammarlens}[title={What you've built}]
One more word for this chapter.
\end{grammarlens}

\subsection*{Your turn}
\begin{itemize}
\raggedright
  \item \emph{Say it:} \emph{der Mensch}
  \item \emph{Say it:} \emph{der Mensch}, once more
  \item \emph{Say it:} say \emph{der Mensch}
  \item \emph{Recall:} say the German for a granddaughter, then the German for a baby, then say \emph{der Mensch} again
\end{itemize}

\begin{tcolorbox}[breakable,colback=teal!4,colframe=teal!35!black,title={Before you move on}]
What does der Mensch mean? (\textquotedblleft{}A human being\textquotedblright{}.) Say it once more.
\end{tcolorbox}
